\newpage
\section*{ABSTRACT}  
\phantomsection
\addcontentsline{toc}{section}{ABSTRACT}

The rapid advancement of artificial intelligence based text-to-speech, voice conversion, and voice cloning technologies has made synthetic speech increasingly realistic and difficult to distinguish from genuine human speech, creating serious risks for VoIP communication, voice authentication, online meetings, call center security, fraud prevention, and digital trust. This report presents the design, implementation, and evaluation of a real-time Audio Deepfake Detection system for identifying AI-generated speech in degraded Voice over Internet Protocol and real-time communication environments. Unlike conventional approaches that mainly evaluate detection models on clean offline recordings, the proposed project focuses on receiver-side reconstructed audio streams affected by practical communication processes such as codec compression, packetization, network delay, jitter, packet loss, packet loss concealment, decoding, buffering, and audio enhancement. The system processes short streaming audio chunks, applies preprocessing and feature extraction using MFCC-based acoustic descriptors and self-supervised speech embeddings such as WavLM, and classifies each segment as real or synthetic using both deep learning architectures (including AASIST and NeXt-TDNN (7.36\% ERR)) and classical machine learning classifiers (including XGBoost and SVM (7.69\% ERR)). Temporal aggregation will be used to stabilize chunk-level predictions and generate confidence scores suitable for real-time alerts, monitoring, and security enforcement. The detector is proposed as an external API-based inference service that can be integrated with RTC or VoIP platforms through client-side or server-side deployment without modifying the core communication stack. As of this progress report, eight baseline detection models spanning both families have been implemented and evaluated on the ASVspoof5 benchmark under clean-audio conditions; integration of the RTC-aware transmission pipeline and the corresponding robustness evaluation under codec compression and packet loss are planned for the final report, alongside assessment of the system's latency, scalability, and deployment feasibility.

\vspace{0.5cm}

\textit{Keywords: Audio Deepfake Detection, Codec Compression, Machine Learning, Packet Loss, Real-Time Communication, Synthetic Speech, Voice over Internet Protocol}